\chapter{Аудит пакета якорей эйнштейновского отклика}
\label{ch:v10-einstein-anchor-package-audit}

\section{Четыре независимых дефицита}

Условная формула
\begin{equation}
 \kappa^2=\frac{24\pi G\Theta\log2}{c^2v_{\rm cell}}
 \label{eq:v10-anchor-package-scale}
\end{equation}
становится физическим выводом только после независимого происхождения
четырёх объектов: ньютоновской постоянной $G$, энергии на единицу энтропии
$\Theta$, абсолютного объёма $v_{\rm cell}$ и вакуумоподобного тензора
энергии сквозного потока.

Для каждого объекта проверены четыре кандидата по пяти одинаковым условиям:
правильный тип и размерность, внутренняя доступность, независимость от
$\kappa$, происхождение из общего родителя и отсутствие круговой подстановки.

\section{Результаты четырёх аудитов}

Для $G$ проверены индуцированный эйнштейновский коэффициент, коэффициент
спектрального действия, импортированная ньютоновская постоянная и решение
$G$ из целевой формулы. Матрица имеет ранг $3$, максимум $3/5$, проходов
$0/4$.

Для $\Theta$ проверены энергия КМС, концевая спектральная щель, внешняя
температура Ландауэра и безразмерное вакуумное действие. Ранг равен $4$,
максимум $4/5$, проходов $0/4$. Ближе всего концевая щель, но её абсолютная
величина не выбрана общим родителем.

Для $v_{\rm cell}$ проверены собственный объём, спектральная счётная мера,
топологический квант и космологический радиус. Ранг равен $2$, максимум
$4/5$, проходов $0/4$: все внутренние варианты сохраняют свободный
размерный множитель.

Для тензора потока проверены скаляр производства энтропии, изотропный
вакуумный анзац, метрическая вариация текущего родителя и тензор отклика
Келдыша. Ранг равен $3$, максимум $4/5$, проходов $0/4$. Первые два варианта
не выводят физический тензор, а последние два ещё не получены вариацией
существующего общего действия.

Объединённая матрица $16\times5$ имеет полный критериальный ранг $5$.
Следовательно, нулевой результат не вызван вырожденным аудитом: все пять
требований различаются, но ни одна строка не выполняет их одновременно.

\section{Пакетный запрет}

Четыре компонента входят в эйнштейновский мост независимо, поэтому их
матрица зависимости равна $I_4$ и имеет ранг $4$. Текущий вектор полной
доступности равен
\begin{equation}
 (0,0,0,0)^T.
 \label{eq:v10-anchor-package-zero-vector}
\end{equation}

Одна формула \eqref{eq:v10-anchor-package-scale} на переменных
$(\kappa,G,\Theta,v_{\rm cell})$ имеет ранг $1$ и ядро размерности $3$.
Если три размерных якоря и тензорное отображение были бы выведены
независимо, полный ранг стал бы равен $4$. Однако комбинирование нескольких
частичных кандидатов не помогает: каждый из них наследует хотя бы одну
невыведенную масштабную или родительскую зависимость.

\begin{theorem}[Пакет якорей действительно четырёхкомпонентен]
Ньютоновская постоянная, энергия энтропии, объём ячейки и тензор энергии
потока являются независимыми требованиями эйнштейновского замыкания. Аудит
шестнадцати кандидатов покрывает все пять критериев, но даёт ноль полных
кандидатов и ноль полностью закрытых компонентов пакета.
\end{theorem}

\begin{nogocriterion}[Запрет сборки масштаба из частичных якорей]
Нельзя объявлять формулу проводимости выведенной, подставляя одновременно
условный $G$, часовую температуру, свободный объём и постулированный
изотропный тензор. Совпадение размерностей такого произведения не устраняет
происхождение ни одного из его свободных множителей.
\end{nogocriterion}

\begin{protocol}[Статус аудита пакета]\begin{enumerate}
\item проверенных кандидатов: \textbf{$16$};
\item полный ранг критериев: \textbf{$5$};
\item индивидуальных проходов: \textbf{$0/16$};
\item закрытых компонентов пакета: \textbf{$0/4$};
\item абсолютная проводимость: \textbf{$0/1$};
\item следующий гейт: \textbf{родитель индуцированной ньютоновской
постоянной}.
\end{enumerate}\end{protocol}

Машинный сертификат сохранён в
\path{s2t_v10_cell_birth_four_volume_einstein_response_anchor_package_candidate_audit_gate_results.json}.